
%=========================================================
% PARTIE 2.a : ETUDE DE L'ARCHITECTURE DIFFWAVE
% ! DEPRECATED : UTILISER PARTIE2_AV2
%=========================================================

\section{Étude de l'architecture DNN DiffWave}

\subsection{Présentation générale de DiffWave}

DiffWave (\textit{Diffusion Model for Audio Synthesis}) est un
modèle génératif profond basé sur les modèles probabilistes
de diffusion. Il a été proposé par Zhifeng Kong et ses
collaborateurs dans l'article intitulé
\textit{DiffWave: A Versatile Diffusion Model for Audio Synthesis},
publié à la conférence ICLR en 2021.

L'objectif principal de DiffWave est de générer des formes
d'onde audio de haute qualité à partir d'un bruit gaussien.
Contrairement aux modèles génératifs autorégressifs classiques,
DiffWave utilise un processus progressif de débruitage pour
transformer un signal initialement aléatoire en un signal audio
structuré.

Le modèle repose sur deux processus complémentaires :

\begin{itemize}
    \item \textbf{Le processus de diffusion (Forward Process)} :
    il consiste à ajouter progressivement du bruit gaussien
    à un signal audio réel jusqu'à obtenir un signal proche
    d'un bruit blanc.

    \item \textbf{Le processus inverse (Reverse Process)} :
    il consiste à apprendre à supprimer progressivement
    le bruit afin de reconstruire un signal audio cohérent.
\end{itemize}

Le processus inverse est appris à l'aide d'un réseau de
neurones profond qui estime le bruit présent dans le signal
à chaque étape de diffusion.

DiffWave est un modèle non autorégressif : il ne génère pas
les échantillons audio un par un, mais traite une forme
d'onde complète à chaque étape du processus de débruitage.

\subsection{Applications de DiffWave}

L'article original présente trois principales applications
de DiffWave dans le domaine de la génération audio.

\subsubsection{Vocodage neuronal (Neural Vocoding)}

Le vocodage neuronal consiste à transformer une représentation
spectrale de la parole en une forme d'onde audio directement
écoutable.

Dans cette application, DiffWave reçoit comme entrée un
spectrogramme Mel, généralement produit par un système
de synthèse vocale à partir d'un texte.

Le spectrogramme Mel fournit des informations temporelles
et fréquentielles sur le signal de parole. Il sert de
conditionnement au modèle afin de guider la génération
de la forme d'onde.

\textbf{Entrée :}
\begin{itemize}
    \item Un spectrogramme Mel correspondant à la parole
    à synthétiser ;
    \item Un signal initial constitué de bruit gaussien
    pendant le processus de génération.
\end{itemize}

\textbf{Sortie :}
\begin{itemize}
    \item Une forme d'onde audio reconstruite, représentant
    le signal de parole correspondant au spectrogramme Mel.
\end{itemize}

Cette application permet notamment de produire une parole
naturelle à partir des caractéristiques spectrales fournies
par un système de synthèse vocale.

\subsubsection{Génération conditionnée par une classe
(Class-Conditional Generation)}

DiffWave peut également générer des signaux audio en
fonction d'une information de classe fournie au modèle.

Cette information peut représenter, par exemple, une
identité de locuteur ou une catégorie de signal.

Le modèle reçoit un bruit gaussien ainsi qu'une étiquette
qui indique le type de signal à générer. Cette information
est utilisée pour orienter le processus de génération.

\textbf{Entrée :}
\begin{itemize}
    \item Un bruit gaussien initial ;
    \item Une étiquette de classe, par exemple un identifiant
    de locuteur.
\end{itemize}

\textbf{Sortie :}
\begin{itemize}
    \item Une forme d'onde audio appartenant à la classe
    spécifiée.
\end{itemize}

Le conditionnement permet ainsi de contrôler certaines
caractéristiques du signal généré.

\subsubsection{Génération non conditionnée
(Unconditional Generation)}

Dans cette application, DiffWave génère un signal audio
sans recevoir d'information conditionnelle externe.

Le modèle commence directement avec un bruit gaussien
aléatoire et applique le processus inverse de diffusion
pour produire une forme d'onde audio.

\textbf{Entrée :}
\begin{itemize}
    \item Un bruit gaussien aléatoire.
\end{itemize}

\textbf{Sortie :}
\begin{itemize}
    \item Une forme d'onde audio générée par le modèle,
    sans conditionnement externe.
\end{itemize}

Cette approche permet au modèle d'apprendre la distribution
des signaux audio présents dans sa base d'apprentissage
et de générer de nouveaux exemples présentant des
caractéristiques similaires.

\subsubsection{Synthèse des applications}

Le tableau suivant résume les trois applications étudiées
dans l'article original :

\begin{table}[H]
\centering
\renewcommand{\arraystretch}{1.5}
\begin{tabular}{|p{3.2cm}|p{5cm}|p{5cm}|}
\hline
\textbf{Application} &
\textbf{Entrée du modèle} &
\textbf{Sortie du modèle} \\
\hline

Vocodage neuronal &
Bruit gaussien et spectrogramme Mel &
Forme d'onde audio correspondant au spectrogramme Mel \\
\hline

Génération conditionnée &
Bruit gaussien et étiquette de classe &
Forme d'onde correspondant à la classe spécifiée \\
\hline

Génération non conditionnée &
Bruit gaussien uniquement &
Nouvelle forme d'onde audio générée sans conditionnement \\
\hline
\end{tabular}
\caption{Principales applications de DiffWave et leurs
entrées et sorties.}
\label{tab:diffwave_applications}
\end{table}

\subsection{Principe de l'approche d'apprentissage de DiffWave}

L'apprentissage de DiffWave repose sur le principe des
modèles probabilistes de diffusion. Il s'effectue en deux
phases principales : un processus direct qui ajoute du bruit
aux données réelles et un processus inverse qui apprend
à reconstruire ces données.

\subsubsection{Processus direct : diffusion (Forward Process)}

Le processus de diffusion consiste à ajouter progressivement
du bruit gaussien à un signal audio réel.

On note le signal audio original par :

\[
x_0 \in \mathbb{R}^{L}
\]

où $L$ représente le nombre d'échantillons du signal.

À chaque étape $t$, un bruit gaussien est ajouté au signal
précédent. Cette opération est définie par :

\[
q(x_t \mid x_{t-1}) =
\mathcal{N}\left(
x_t;
\sqrt{1-\beta_t}\,x_{t-1},
\beta_t I
\right)
\]

où :

\begin{itemize}
    \item $x_t$ représente le signal après l'étape de diffusion
    $t$ ;
    \item $\beta_t$ est un paramètre qui contrôle la quantité
    de bruit ajoutée à cette étape ;
    \item $I$ représente la matrice identité ;
    \item $\mathcal{N}$ désigne une distribution gaussienne.
\end{itemize}

Les valeurs de $\beta_t$ suivent un calendrier de variance
(\textit{variance schedule}) défini avant l'apprentissage.

En répétant cette opération, le signal original perd
progressivement ses caractéristiques et devient de plus
en plus bruité.

Après un nombre suffisamment élevé d'étapes $T$, le signal
$x_T$ se rapproche d'un bruit gaussien standard :

\[
x_T \sim \mathcal{N}(0,I).
\]

Une propriété importante de ce processus est qu'il est
possible de calculer directement le signal bruité à une
étape quelconque $t$, sans exécuter toutes les étapes
intermédiaires.

En posant :

\[
\alpha_t=1-\beta_t
\]

et

\[
\overline{\alpha}_t=
\prod_{s=1}^{t}\alpha_s,
\]

on obtient :

\[
\boxed{
x_t =
\sqrt{\overline{\alpha}_t}\,x_0
+
\sqrt{1-\overline{\alpha}_t}\,\epsilon
}
\]

avec :

\[
\epsilon\sim\mathcal{N}(0,I).
\]

Cette équation permet de générer directement une version
bruitée du signal original pour une étape de diffusion
choisie aléatoirement.

\subsubsection{Processus inverse : débruitage
(Reverse Process)}

Le processus inverse représente la phase de génération
du signal audio.

Son objectif est de partir d'un bruit gaussien $x_T$
et de reconstruire progressivement un signal audio
cohérent en supprimant le bruit ajouté pendant le
processus direct.

Le modèle apprend une distribution conditionnelle
de transition :

\[
p_{\theta}(x_{t-1}\mid x_t)
\]

qui permet de passer du signal bruité $x_t$ au signal
moins bruité $x_{t-1}$.

Cette distribution est modélisée par :

\[
p_{\theta}(x_{t-1}\mid x_t)
=
\mathcal{N}
\left(
x_{t-1};
\mu_{\theta}(x_t,t),
\sigma_t^2 I
\right)
\]

où $\mu_{\theta}$ représente la moyenne prédite par
le modèle et $\sigma_t$ son écart-type.

Le réseau de neurones $\epsilon_{\theta}$ apprend à
estimer le bruit présent dans le signal reçu.

La moyenne de la distribution inverse est alors
paramétrée par :

\[
\boxed{
\mu_{\theta}(x_t,t)=
\frac{1}{\sqrt{\alpha_t}}
\left(
x_t-
\frac{\beta_t}
{\sqrt{1-\overline{\alpha}_t}}
\epsilon_{\theta}(x_t,t)
\right)
}
\]

À chaque étape, le modèle utilise cette estimation
pour reconstruire un signal moins bruité.

Le processus est répété de $t=T$ jusqu'à $t=1$.
Le résultat final est le signal généré $x_0$.

\subsubsection{Illustration des deux processus}

Le fonctionnement général de DiffWave peut être résumé
par le schéma suivant :

\begin{center}
\fbox{
\begin{minipage}{0.88\textwidth}
\centering

\textbf{Processus direct : ajout progressif de bruit}

\vspace{0.3cm}

Signal audio réel $x_0$

$\Downarrow$

Signal légèrement bruité $x_1$

$\Downarrow$

Signal de plus en plus bruité $x_2,\ldots,x_{T-1}$

$\Downarrow$

Bruit gaussien $x_T$

\vspace{0.4cm}
\hrule
\vspace{0.4cm}

\textbf{Processus inverse : génération par débruitage}

\vspace{0.3cm}

Bruit gaussien $x_T$

$\Downarrow$

Estimation et suppression progressive du bruit

$\Downarrow$

Signaux intermédiaires $x_{T-1},\ldots,x_1$

$\Downarrow$

Signal audio généré $x_0$

\end{minipage}
}
\end{center}

Le processus direct est utilisé pour construire des
exemples bruités pendant l'apprentissage, tandis que
le processus inverse est appris par le réseau de neurones
et utilisé pour générer de nouveaux signaux audio.

\subsection{Apprentissage du réseau de neurones}

Pendant l'apprentissage, le modèle reçoit un signal
audio original $x_0$, une étape de diffusion $t$ et
un bruit gaussien aléatoire $\epsilon$.

À partir de ces éléments, il construit un signal bruité :

\[
x_t =
\sqrt{\overline{\alpha}_t}\,x_0
+
\sqrt{1-\overline{\alpha}_t}\,\epsilon.
\]

Le réseau de neurones reçoit alors $x_t$ et l'indice
de diffusion $t$. Son objectif est de prédire le bruit
$\epsilon$ qui a été ajouté au signal.

La prédiction est notée :

\[
\epsilon_{\theta}(x_t,t).
\]

La fonction de coût utilisée par DiffWave est une erreur
quadratique entre le bruit réel et le bruit estimé :

\[
\boxed{
\mathcal{L}(\theta)=
\mathbb{E}_{x_0,\epsilon,t}
\left[
\left\|
\epsilon-
\epsilon_{\theta}
\left(
\sqrt{\overline{\alpha}_t}x_0+
\sqrt{1-\overline{\alpha}_t}\epsilon,
t
\right)
\right\|_2^2
\right]
}
\]

L'objectif de l'apprentissage est de minimiser cette
fonction de coût en ajustant les paramètres $\theta$
du réseau.

En pratique, l'étape de diffusion $t$ est choisie
aléatoirement parmi les étapes possibles. Le modèle
apprend ainsi à estimer le bruit pour différents niveaux
de dégradation du signal.

Cette méthode permet d'entraîner le réseau à effectuer
le débruitage à différents niveaux, sans devoir lui
fournir séparément un réseau discriminateur ou un
encodeur supplémentaire.

\subsection{Conditionnement dans DiffWave}

Le conditionnement consiste à fournir au modèle une
information supplémentaire afin de contrôler ou
d'orienter le signal audio généré.

Dans DiffWave, le réseau peut fonctionner sans
conditionnement ou avec une information conditionnelle.

On distingue principalement deux types de conditionnement.

\subsubsection{Conditionnement local (Local Conditioning)}

Le conditionnement local fournit une information qui
varie dans le temps et qui est alignée avec les
échantillons du signal audio.

Dans le cas du vocodage neuronal, cette information
correspond au spectrogramme Mel.

Le spectrogramme Mel est d'abord mis à l'échelle
temporelle afin d'avoir une longueur compatible avec
celle de la forme d'onde.

Il est ensuite injecté dans les couches du réseau
afin de guider la génération du signal audio.

Ainsi, le modèle ne génère pas une parole quelconque :
il produit une forme d'onde qui respecte les
caractéristiques temporelles et spectrales du
spectrogramme fourni.

\subsubsection{Conditionnement global (Global Conditioning)}

Le conditionnement global fournit une information
qui caractérise l'ensemble du signal plutôt que chaque
instant individuellement.

Il peut s'agir, par exemple, d'un identifiant de
locuteur ou d'une étiquette de classe.

Cette information est transformée en représentation
numérique à l'aide d'une couche d'embedding, puis
transmise au réseau.

Elle permet d'orienter la génération vers une catégorie
ou un locuteur particulier.

\subsubsection{Différence entre diffusion et conditionnement}

Il est important de distinguer les deux notions :

\begin{itemize}
    \item La \textbf{diffusion} décrit le mécanisme
    mathématique d'ajout progressif de bruit puis de
    suppression de ce bruit.

    \item Le \textbf{conditionnement} représente
    l'information supplémentaire qui guide la génération
    du signal vers un résultat souhaité.
\end{itemize}

La diffusion répond donc à la question :
\textit{comment le signal est-il progressivement généré ?}

Le conditionnement répond à la question :
\textit{quel type de signal le modèle doit-il générer ?}

\subsection{Différences entre DiffWave et les vocodeurs
classiques tels que WaveNet}

WaveNet est un modèle génératif autorégressif qui
génère les échantillons audio de manière séquentielle.
Chaque nouvel échantillon dépend des échantillons
précédemment générés.

DiffWave adopte une approche différente : il utilise
un processus de diffusion probabiliste pour générer
une forme d'onde à partir d'un bruit initial.

Bien que son architecture utilise également des
convolutions dilatées et des connexions résiduelles,
son mécanisme de génération n'est pas autorégressif.

Les principales différences sont présentées dans
le tableau suivant.

\begin{table}[H]
\centering
\renewcommand{\arraystretch}{1.5}
\begin{tabular}{|p{3.1cm}|p{5.1cm}|p{5.1cm}|}
\hline
\textbf{Critère} &
\textbf{WaveNet} &
\textbf{DiffWave} \\
\hline

Principe de génération &
Génération autorégressive &
Génération par diffusion probabiliste \\
\hline

Ordre de génération &
Échantillon par échantillon &
Signal complet traité à chaque étape inverse \\
\hline

Dépendance temporelle &
Convolutions causales &
Convolutions dilatées bidirectionnelles \\
\hline

Entrée initiale &
Contexte des échantillons précédents &
Bruit gaussien initial \\
\hline

Mécanisme principal &
Prédiction du prochain échantillon &
Estimation et suppression progressive du bruit \\
\hline

Conditionnement &
Caractéristiques locales ou linguistiques &
Conditionnement local ou global possible \\
\hline

Parallélisation &
Génération séquentielle &
Calcul parallèle des échantillons à chaque étape \\
\hline

Temps de génération &
Peut être important pour les longues formes d'onde &
Réduit grâce à un nombre limité d'étapes de diffusion \\
\hline
\end{tabular}
\caption{Comparaison entre les approches de génération
audio de WaveNet et DiffWave.}
\label{tab:diffwave_wavenet}
\end{table}

Dans l'article original, les auteurs rapportent que
DiffWave atteint un score MOS de 4.44, contre 4.43
pour le vocodeur WaveNet utilisé dans leur comparaison.
Le score MOS (\textit{Mean Opinion Score}) est une
mesure subjective de la qualité perçue de la parole.

Les auteurs montrent également que DiffWave peut
synthétiser les formes d'onde beaucoup plus rapidement
que WaveNet dans leur configuration expérimentale.

Il faut toutefois noter que ces résultats dépendent
des configurations expérimentales et ne signifient
pas que DiffWave est systématiquement plus rapide
que tous les vocodeurs existants.

\subsection{Conclusion}

DiffWave est un modèle génératif audio basé sur les
modèles probabilistes de diffusion. Son fonctionnement
repose sur un processus direct d'ajout de bruit et un
processus inverse de débruitage appris par un réseau
de neurones profond.

Sa capacité à fonctionner avec ou sans conditionnement
lui permet de prendre en charge plusieurs applications,
notamment le vocodage neuronal, la génération
conditionnée par une classe et la génération non
conditionnée.

Contrairement à WaveNet, qui génère les échantillons
de manière autorégressive et séquentielle, DiffWave
utilise un processus de génération non autorégressif
reposant sur plusieurs étapes successives de débruitage.

Cette approche permet de produire des formes d'onde
audio de haute qualité tout en offrant une alternative
intéressante aux vocodeurs autorégressifs classiques.

%=========================================================
% REFERENCE BIBLIOGRAPHIQUE
%=========================================================

% \begin{thebibliography}{9}

% \bibitem{kong2021diffwave}
% Z. Kong, W. Ping, J. Huang, K. Zhao et B. Catanzaro,
% \textit{DiffWave: A Versatile Diffusion Model for Audio Synthesis},
% International Conference on Learning Representations (ICLR), 2021.

% \end{thebibliography}
